\section{Related Work}
\label{sec:related_work}

\paragraph{Memory in Transformers.}
Transformers implement in-context memory primarily through self-attention, where
each token can retrieve information from earlier positions by computing
content-dependent weights over key--value pairs \cite{vaswani2017attention}.
This mechanism is often viewed as a form of differentiable, content-addressable
memory: rather than compressing the past into a single recurrent state, the
model maintains access to a growing set of past representations, enabling
flexible retrieval and composition over the entire context window. A large body
of work has therefore focused on extending the effective context length or
augmenting the model with additional memory beyond the fixed window. Notable
approaches include recurrence across segments (Transformer-XL), which reuses
previous hidden states to model dependencies beyond a fixed-length context
\cite{dai2019transformerxl}, compressing older memories to preserve long-range
information under a bounded compute budget \cite{rae2019compressive}, and
retrieval-augmented language modeling, which explicitly consults an external
datastore during generation \cite{borgeaud2022retro}. These methods primarily
target \emph{memory availability} (how much information can be brought into the
computation) and \emph{access} (how that information can be retrieved when
needed).

%\paragraph{Sub-quadratic attention and memory.}
%Since standard attention scales quadratically with sequence length, efficient
%Transformer variants aim to preserve attention-like retrieval while reducing
%time and/or memory. Dense approximations typically exploit structure in the
%attention matrix, such as low-rank projections (Linformer)
%\cite{wang2020linformer} or kernel-based approximations with random features
%(Performer) \cite{choromanski2021performer}. Other approaches use approximate
%nearest-neighbor search or hashing to restrict attention computation (Reformer)
%\cite{kitaev2020reformer}. These models are frequently evaluated on long-context
%benchmarks designed to stress memory and long-range dependency modeling, such as
%Long Range Arena (LRA) \cite{tay2020lra}. A recurring theme across this
%literature is that the \emph{form} of the attention approximation matters:
%changing the mechanism used for retrieval (e.g., low-rank vs.\ kernel vs.\
%sparsity) can alter which long-range dependencies remain accessible, thereby
%affecting apparent memory performance even when the nominal context length is
%unchanged.

\paragraph{Attention-free sub-quadratic sequence models.}
A complementary line of work achieves sub-quadratic (often linear-time)
inference by replacing attention with recurrent or convolutional updates that
compress the past into a fixed-size state. State space models (SSMs) provide a
particularly unified view: they maintain an internal state updated online and
produce outputs via a readout from that state. Foundational work on HiPPO
introduced principled mechanisms for online compression of history via
projections \cite{gu2020hippo}. Building on these ideas, structured
parameterizations such as S4 made SSMs computationally practical and competitive
on long-range tasks \cite{gu2021s4}, and subsequent simplifications like DSS
showed that diagonal state parameterizations can retain much of S4's
effectiveness \cite{gupta2022dss}. Modern language-model-scale architectures
(e.g., selective SSMs and gated linear attention variants) can often be
interpreted as instances of this broader paradigm: the model stores information
in a recurrent state and must \emph{query} or \emph{read} from it during
generation. This raises a distinct memory question from the Transformer setting:
when the past is compressed into a fixed-size state, the limiting factor may be
less about what can be stored in principle and more about which stored
information is practically retrievable given the model's readout mechanism.

\paragraph{Entity binding as an in-context memory primitive.}
Entity--attribute binding tasks probe a core capability underlying factual
recall and multi-entity reasoning: given a context describing several entities
and their properties, the model must correctly bind each property to its
corresponding entity at query time. Recent work has used controlled settings to
study how binding is represented internally; for example, Feng and Steinhardt
identify a ``binding ID'' mechanism and use causal interventions to argue that
models attach binding-related vectors to entity and attribute positions
\cite{feng2024binding}. Such findings motivate synthetic entity-binding
evaluations as a targeted lens on in-context memory that avoids confounds from
pretraining knowledge. Our entity-binding datasets are in this spirit, but we
focus on sub-quadratic models whose memory is mediated by a recurrent state:
this shifts the emphasis from \emph{Can tokens attend to the right binding} to
\emph{which bindings are present in the compressed state and how the model
queries them}. In particular, by comparing what can be decoded from the state to
what the model actually retrieves during generation, we align entity binding
with a concrete hypothesis about the failure mode in sub-quadratic
architectures: a potential mismatch between \emph{stored} information and the
model's \emph{query} mechanism.

\paragraph{Interpretability tools for studying memory mechanisms.}
Our work sits within mechanistic interpretability, which aims to localize and
characterize internal computations responsible for model behavior. Probing
methods train auxiliary predictors to decode information from hidden
activations; recent 'lens' techniques decode intermediate representations into
output distributions, providing a token-level view of how predictions evolve
across depth \cite{belrose2023tunedlens}. Beyond correlational probes, causal
interventions are commonly used to establish necessity and sufficiency of
components. Meng et al.\ develop causal tracing interventions to localize
computations supporting factual recall and show that specific mechanisms can be
edited (ROME) \cite{meng2022rome}. Activation patching has become a standard causal tool for circuit-level analysis, and
recent work systematizes methodological choices and metrics for patching-based
localization \cite{zhang2023patching}. Broader frameworks for circuit-based
reasoning about Transformers provide a language for decomposing computations
into interpretable parts \cite{elhage2021framework}. We adopt these
interpretability principles in the sub-quadratic setting by (i) designing probes
that mirror the architecture's native readout direction and (ii) using
counterfactual patching on internal state representations to test whether
probe-identified components have a causal effect on retrieval.
